\section{Norm comparisons from an auxiliary pair}
\label{sec:comparators}

The metric deformation of Section~\ref{sec:transport} provides a vector
whose physical energy can be small.  We now compare the norm defining that
vector in two ways.  Projecting the target and an auxiliary system onto an
entangled pair gives an upper bound on the negative logarithm of the norm.
A lower bound follows because most physical copies of a low-energy vector
are in the ground state.  The auxiliary systems retain large representation
labels in this second comparison.  Their contribution is what makes the two
bounds useful.

We formulate the comparison for a finite collection of partitions.  No
traversal procedure or choice of geometric scales is needed in this section.
All physical entropies below are in a fixed normalized pure vector
$\widetilde\Omega\in\mathcal H_{\mathrm{phys}}$, unless another state is
specified.  Suppose that a positive Hamiltonian $\widetilde H$ has this
unique ground vector, ground energy $\widetilde E_0$, and gap
$\widetilde g>0$.  Fix a physical subsystem $X$.  A status consists of $K$
partitions
\[
 \mathcal H_{\mathrm{phys}}=\mathcal H_X\otimes
 \mathcal H_{U_g}\otimes\mathcal H_{Y_g}\otimes\mathcal H_{V_g},
 \qquad 1\le g\le K.
\]
There are finitely many allowed statuses.  We assume that, for every allowed
choice of the two statuses and every $g<g'$, their physical sets satisfy
\begin{equation}
 XU_gY_g\subseteq XU_{g'}.
 \label{comparator:nesting}
\end{equation}
As usual, juxtaposition joins disjoint tensor factors.  The inclusion in
\eqref{comparator:nesting} compares sets of physical sites.  It ensures the
commutation of the different bands' metric algebras after complementary
labels are identified on the symmetric space.

Here are the entropy data used in the comparisons.  Write
$S_X=S(X)$.  Choose a nonempty set $E$ of positive Schmidt indices of $X$
such that, for some $w\ge0$,
\begin{equation}
 e^{-S_X-w}\le p_i\le e^{-S_X+w}\quad(i\in E),
 \qquad z=\sum_{i\in E}p_i>0.
 \label{comparator:typical-set}
\end{equation}
Put $d_*=|E|$, $p_i'=p_i/z$, and $S'= -\sum_{i\in E}p_i'\log p_i'$.
For the later choice of a scale parameter $n\ge2$,
Lemma~\ref{lem:tail} will supply $w=n^{3/5}$ and $z\ge1-n^{-100}$.
Only \eqref{comparator:typical-set} is needed here.  Directly from
\eqref{comparator:typical-set},
\begin{equation}
 |\log d_*-S_X|\le w+|\log z|,
 \qquad |S'-S_X|\le w+|\log z|.
 \label{comparator:typical-entropies}
\end{equation}

For the rough upper bound, assume there is a finite nested collection
$\mathcal R$ of physical subsystems disjoint from $X$.  For every status
and band, each of $U_g$ and $U_gY_g$ has a matched member of $\mathcal R$
which contains it or is contained in it.  Assume
\begin{equation}
 S(B)\le B_{\mathrm{sh}}\quad(B\in\mathcal R),
 \qquad 1+\log\dim\mathcal H_{B\triangle Q}\le B_{\mathrm{exc}}
 \label{comparator:prefix-data}
\end{equation}
for every such matched pair $(B,Q)$, where $Q=U_g$ or $U_gY_g$.
The symmetric difference here is simply the difference between the two
nested sets.  The added $1$ in $B_{\mathrm{exc}}$ absorbs a fixed margin
in a spectral cutoff.

The sharp bounds additionally use the following marginal moment assumption.
For every set
\[
 Q\in\{XU_g,V_g,Y_g,U_g,U_gY_g:\text{allowed statuses and bands}\},
\]
let $\rho_Q$ be the marginal of $\widetilde\Omega$.  There are constants
$C_0,c_0>0$ and $\mathcal B\ge1$ such that
\begin{equation}
 \log\Tr\rho_Q^{\,1-u}
 \le uS(Q)+C_0\mathcal B u^2,
 \qquad |u|\le c_0/\sqrt{\mathcal B}.
 \label{comparator:mgf-data}
\end{equation}
Only positive eigenvalues are used in this expression.  We retain the
parameter $\mathcal B$ from the marginal estimate, including any
logarithmic factors it contains.

Append systems $C,R$ of dimension $d_*$, with bases indexed by $E$, and put
\[
 \mathrm{Bell}_{CR}=d_*^{-1/2}\sum_{i\in E}|i\rangle_C|i\rangle_R.
\]
Use the common-dimension Schur metrics of Section~\ref{sec:replicas} on
\[
 \mathcal V=\mathcal H_{\mathrm{phys}}\otimes\mathcal H_C
       \otimes\mathcal H_R,
 \qquad \mathcal S_k=\Sym^k\mathcal V.
\]
In particular the same scalar function of a partition label defines every
$W_{Q,k}$, with labels padded to $\dim\mathcal V$.  Suppress $k$ when it
is unambiguous.  Fix $0<a\le a_0$, where $a_0>0$ is a sufficiently small
absolute constant, set $t=a/2$, $s=1/4$, and $W=aK$.  At a terminal leaf
$j$ of any finite binary mean tree, put
\begin{equation}
 P_{j,g}=CXU_{j,g},\qquad F_{j,g}=RV_{j,g},\qquad
 A_{j,g}=(W_{P_{j,g}}^{-1}W_{F_{j,g}}^{-1}W_{Y_{j,g}})^2,
 \qquad A_j=\prod_{g=1}^K A_{j,g}.
 \label{comparator:metrics}
\end{equation}
Thus $X,U_{j,g},Y_{j,g},V_{j,g}$ are the physical factors, while
$P_{j,g},Y_{j,g},F_{j,g}$ form the full partition used by the metric.
The target is $X$. The auxiliary factors $C,R$ are initialized in
the Bell pair above, and the same factors are used in every band.
Only the physical factors contribute to the Hamiltonian.

Let $w_j$ be the terminal path probabilities, let $\mathcal M$ be the
mean-tree root with leaves $A_j$, and let $\mathcal M_g$ use the same
tree with leaves $A_{j,g}$.  Then
$\mathcal M=\prod_g\mathcal M_g$, with commuting factors.  We write
\[
 \langle f\rangle_{\mathrm{av}}=
       \frac1K\sum_{g=1}^K\sum_jw_j f_{j,g}
\]
for the full terminal-path and band average.  In particular it includes
both old and new terminal leaves when the tree contains an interpolation.

\begin{proposition}[Upper and lower norm comparisons]
\label{prop:comparators}
Fix the positive gapped Hamiltonian and physical target $X$ above,
a finite collection of status partitions satisfying
\eqref{comparator:nesting}, a typical Schmidt set satisfying
\eqref{comparator:typical-set}, and a mean tree with metrics
\eqref{comparator:metrics}. There is a sequence of common auxiliary
labels $\lambda_k$ and nonzero vectors
\begin{equation}
 \mathrm{pre}_k=\pi_R^{\lambda_k}
  (\widetilde\Omega\otimes\mathrm{Bell}_{CR})^{\otimes k}
       \in\mathcal S_k,
 \qquad \|\mathrm{pre}_k\|\le1,
 \label{comparator:prevector}
\end{equation}
Here $\pi_R^{\lambda_k}$ is the Schur-label projection on the $R$ copies.
These vectors are physical mean-energy eigenvectors of eigenvalue
$\widetilde E_0$ and have label $\lambda_k$ on both $C$ and $R$.
For every tree above, set
\[
 N=\|\mathcal M^{-s}\mathrm{pre}_k\|,
 \qquad v=\mathcal M^{-s}\mathrm{pre}_k/N,
 \qquad
 E_{\mathrm{def}}=
 \frac{\langle v,\overline{\widetilde H}v\rangle-\widetilde E_0}
      {\widetilde g},
 \qquad
 \overline{\widetilde H}=\frac1k\sum_{i=1}^k\widetilde H^{(i)}.
\]
\emph{Rough comparisons.} For $0<\tau<1/2$, the upper bound below
uses the prefix data \eqref{comparator:prefix-data}:
\begin{align}
 -\frac{\log N^2}{k}
 &\le 2\log d_*-\log z+
       2sW\left(\frac{2B_{\mathrm{sh}}}{z}+2B_{\mathrm{exc}}\right)
       +o_k(1),
 \label{comparator:rough-upper}\\
 -\frac{\log N^2}{k}
 &\ge 2sW\left[
       2S'-C\left(\tau\log d_*+\frac{h(\tau)}a\right)
       -C\frac{E_{\mathrm{def}}}{\tau}S'\right]-o_k(1).
 \label{comparator:rough-lower}
\end{align}
Here $h$ is binary entropy.  The lower bound does not require
\eqref{comparator:prefix-data} or \eqref{comparator:mgf-data}.

\emph{Sharp comparisons.} If \eqref{comparator:mgf-data} holds and
$a\le c/\sqrt{\mathcal B}$, with $c$ sufficiently small in terms of
$c_0$, define
\begin{align*}
 g_{j,g}^{\mathrm{pre}}
   &=S(XU_{j,g})+S(V_{j,g})-S(Y_{j,g}),\\
 g_{j,g}^{\mathrm{post}}
   &=S(U_{j,g})+S(U_{j,g}Y_{j,g})-S(Y_{j,g}),
 \qquad g_{\max}=\max_{j,g}g_{j,g}^{\mathrm{pre}}.
\end{align*}
Then $g_{j,g}^{\mathrm{pre}}\ge0$, and the sharp comparisons are
\begin{align}
 -\frac{\log N^2}{k}
 &\le 2\log d_*-\log z+
  2sW\left[\langle g^{\mathrm{post}}\rangle_{\mathrm{av}}
          +C\left(a\mathcal B+\frac{|\log z|}{a}\right)\right]
          +o_k(1),
 \label{comparator:sharp-upper}\\
 -\frac{\log N^2}{k}
 &\ge 2sW\left[2S'+(1-\tau)
                  \langle g^{\mathrm{pre}}\rangle_{\mathrm{av}}
       -C\left(\tau\log d_*+\frac{h(\tau)}a+a\mathcal B\right)
       -C\frac{E_{\mathrm{def}}}{\tau}(S'+g_{\max})\right]
       -o_k(1).
 \label{comparator:sharp-lower}
\end{align}
The upper sharp bound does not require \eqref{comparator:prefix-data}.
For every physical partition there is the exact identity
\begin{equation}
 2S_X+g_{j,g}^{\mathrm{pre}}-g_{j,g}^{\mathrm{post}}
       =I(X:V_{j,g})+I(X:Y_{j,g}V_{j,g}).
 \label{comparator:entropy-identity}
\end{equation}

All replica limits are taken with the physical system, auxiliary dimensions,
finite status collection and tree, $a$, and $\tau$ fixed.  The errors are
uniform over terminal probabilities, including their dependence on $k$,
and over the symmetric vectors to which the lower comparison is applied.
Constants or degrees in intermediate polynomial factors in $k$ may depend
on the fixed entire system.  They disappear in $o_k(1)$.  The displayed
constants in the rough bounds are universal; those in the sharp bounds
may also depend on $C_0,c_0$.  None has additional dependence on the
system dimensions.
\end{proposition}

\subsection{Labels selected by the auxiliary pair}

We first record the entropy concentration of the labels that will be used.
For a permutation-invariant density $\rho$ on $k$ copies of a fixed region,
let $L=-\log\rho$ on its support and let $F$ denote the logarithm of its
Specht dimension. Lemma~\ref{lem:schur}(5) gives
\begin{equation}
 F\le L,\qquad [F,L]=0,\qquad
 \Tr\rho\,e^{b(L-F)}\le\poly(k),\quad 0\le b\le1.
 \label{comparator:label-remainder}
\end{equation}
For an iid marginal satisfying \eqref{comparator:mgf-data}, it follows that
\begin{equation}
 \log\E\exp\{\pm u(F_Q-kS(Q))\}
       \le Ck\mathcal B u^2+O(\log(k+1)),
 \qquad 0\le u\le c/\sqrt{\mathcal B}.
 \label{comparator:signed-label-moments}
\end{equation}
For the positive sign use $F_Q\le L$ and the iid moment bound.
For the negative sign apply Cauchy--Schwarz to
$e^{-u(L-kS(Q))}e^{u(L-F_Q)}$, using the moment bound at $-2u$ and
\eqref{comparator:label-remainder} at $2u$.  Shrinking $c$ ensures both
conditions.  Only the $O(\log(k+1))$ term has fixed-system dependence.

Choose Schmidt vectors so that
\[
 \widetilde\Omega=\sum_i\sqrt{p_i}\,|i\rangle_X|w_i\rangle,
 \qquad
 \psi=\sum_{i\in E}\sqrt{p_i'}\,|i\rangle_R|w_i\rangle.
\]
Thus $\psi$ is the normalized typical truncation with $X$ relabeled as
$R$.  Define the vector and projection used to pin the target:
\[
 B_{XC}=\frac1{\sqrt{d_*}}\sum_{i\in E}|i\rangle_X|i\rangle_C,
 \qquad
 \mathrm{post}_k=B_{XC}^{\otimes k}\otimes\psi^{\otimes k},
\]
and let $\mathsf P$ project $(XC)^{\otimes k}$ onto
$B_{XC}^{\otimes k}$.
The auxiliary label must have at least inverse polynomial probability in
$\mathrm{post}_k$.  We therefore select it using $\psi$, then impose
it on the initial auxiliary pair.  Its initial probability need only
be nonzero; the prevector remains unnormalized.

In $\psi^{\otimes k}$, the iid surprisal on $R$ is
$kS'+o(k)$ in probability.  For instance, at fixed system Chebyshev's
inequality gives a window of width $k^{3/4}$.  By
\eqref{comparator:label-remainder}, the difference between that surprisal
and $F_R$ exceeds $k^{3/4}$ with probability at most
$\poly(k)e^{-k^{3/4}/2}$.  The number of labels is polynomial.
Consequently one can select $\lambda_k$ such that
\begin{equation}
 q_k:=\|\pi_R^{\lambda_k}\psi^{\otimes k}\|^2
       \ge\poly(k)^{-1},\qquad
 \log d_{\lambda_k}=kS'+o(k).
 \label{comparator:high-label}
\end{equation}
The same label has positive probability on the initial uniform pair:
every label allowed for $d_*$ occurs in the tensor power of a full-rank
maximally mixed density.  This proves that \eqref{comparator:prevector}
is nonzero.  Its physical part is unchanged by the auxiliary projection,
so it is an exact mean-energy eigenvector.  On the repeated uniform pair,
the $C$ and $R$ labels agree.  Each metric preserves these two labels:
a central permutation operator on an auxiliary system commutes with the
simultaneous permutations on a larger region containing it.  The tree
preserves them as well.

The pinning calculation in one copy, and then in
$k$ copies, gives
\begin{equation}
 \mathsf P\mathrm{pre}_k
       =\left(\frac{\sqrt z}{d_*}\right)^k
                         \pi_R^{\lambda_k}\mathrm{post}_k.
 \label{comparator:bell-pin}
\end{equation}
The projector $\mathsf P$ commutes with all metrics: $XC$ is contained
in every $P_{j,g}$, and the repeated vector $B_{XC}^{\otimes k}$ is
invariant under simultaneous permutations on $XC$.  In
\eqref{comparator:bell-pin} the label projector is on $R$, disjoint from
the pin.  Therefore
\begin{equation}
 N^2\ge\left(\frac z{d_*^2}\right)^k
       \|\mathcal M^{-s}\pi_R^{\lambda_k}\mathrm{post}_k\|^2.
 \label{comparator:pin-norm}
\end{equation}

On the pinned symmetric range, the $P_{j,g}$ label is the $U_{j,g}$
label.  By complementarity the $F_{j,g}$ label is that of $U_{j,g}Y_{j,g}$.
Because the same label function defines all $W_Q$, restriction of the
metric is exactly
\begin{equation}
 \widetilde A_{j,g}
       =(W_{U_{j,g}}^{-1}W_{U_{j,g}Y_{j,g}}^{-1}W_{Y_{j,g}})^2,
 \qquad
 \log\widetilde A_{j,g}
       =a(F_{U_{j,g}}+F_{U_{j,g}Y_{j,g}}-F_{Y_{j,g}})
                         +O(\log(k+1))\id.
 \label{comparator:pinned-metric}
\end{equation}
The final error notation means a bound in operator norm; it does not
assert that the error is scalar.  The restricted tree will be denoted
$\widetilde{\mathcal M}$, with factors
$\widetilde{\mathcal M}_g$.

\subsection{Upper comparisons}

For any physical $B$ disjoint from $X$, tracing the Schmidt truncation
gives
\begin{equation}
 \rho_B=z\rho_{\psi,B}+(1-z)\rho_{\mathrm{rem},B},
 \qquad
 \rho_{\psi,B}\le z^{-1}\rho_B,
 \qquad S(\rho_{\psi,B})\le S(B)/z,
 \label{comparator:post-marginal}
\end{equation}
where the remainder is omitted if $z=1$.  The entropy inequality is
concavity, with the nonnegative remainder entropy discarded.

For the rough comparison, simultaneously impose the cutoffs
\[
 F_B\le k(S(B)/z+1),\qquad B\in\mathcal R,
\]
on $\mathrm{post}_k$.  These spectral projections commute, since
$\mathcal R$ is nested.  Each failure has exponentially small probability
in $k$: under the iid marginal of $\psi$, $F_B$ is at most the
surprisal, whose one-copy law is bounded on its positive spectrum at this
fixed system; exponential Markov and
\eqref{comparator:post-marginal} then apply.  The finite union of failures
is still exponentially small.  All these projections commute with
$\pi_R^{\lambda_k}$.  Their common projection of
$\pi_R^{\lambda_k}\mathrm{post}_k$ therefore has squared norm at least
$\poly(k)^{-1}$ by \eqref{comparator:high-label}.  Normalize it to
$\varphi$; then
\begin{equation}
 |\langle\varphi,\pi_R^{\lambda_k}\mathrm{post}_k\rangle|^2
                         \ge\poly(k)^{-1}.
 \label{comparator:rough-overlap}
\end{equation}

If $Q$ is matched to $B$, the nested label inequalities give
$F_Q\le F_B+k\log\dim\mathcal H_{B\triangle Q}$, in either inclusion
direction.  The two label operators commute.  Thus $\varphi$ belongs to
the spectral range
\begin{equation}
 F_Q\le k(B_{\mathrm{sh}}/z+B_{\mathrm{exc}}),
       \qquad Q=U_{j,g},\ U_{j,g}Y_{j,g}.
 \label{comparator:actual-cutoffs}
\end{equation}
An actual label need only commute with its matched cutoff to obtain this
conclusion; no commutation with every imposed cutoff is being assumed.
Within any one product leaf, all the labels in
\eqref{comparator:pinned-metric} commute by
\eqref{comparator:nesting}.  Dropping the negative labels therefore gives
\[
 \langle\varphi,\prod_g\widetilde A_{j,g}\varphi\rangle
 \le\poly(k)\exp\{2kaK(B_{\mathrm{sh}}/z+B_{\mathrm{exc}})\}.
\]
Equation~\eqref{transport:vector-jensen}, iterated through the tree, gives
the same bound for $\langle\varphi,\widetilde{\mathcal M}\varphi\rangle$.
Since $2s\le1$, another scalar Jensen
inequality bounds the expectation of
$\widetilde{\mathcal M}^{2s}$ by the $2s$ power of its expectation.
Cauchy--Schwarz in \eqref{comparator:rough-overlap}, applied to
$\widetilde{\mathcal M}^{s}\varphi$ and
$\widetilde{\mathcal M}^{-s}\pi_R^{\lambda_k}\mathrm{post}_k$, now proves
\eqref{comparator:rough-upper} using \eqref{comparator:pin-norm}.

For the sharp bound, instead normalize
$\pi_R^{\lambda_k}\mathrm{post}_k$ itself to a vector $\widehat\psi_k$.
Scalar Jensen for the exponential gives
\[
 \|\widetilde{\mathcal M}^{-s}\pi_R^{\lambda_k}\mathrm{post}_k\|^2
 \ge q_k\exp\{-2s\langle\widehat\psi_k,
                    \log\widetilde{\mathcal M}\,\widehat\psi_k\rangle\}.
\]
Split the logarithm into the commuting band factors.  On each band, log
Jensen, commutation with the auxiliary sector, and vector mean Jensen give
\begin{align}
 \langle\widehat\psi_k,
       \log\widetilde{\mathcal M}_g\widehat\psi_k\rangle
 &\le\log\langle\widehat\psi_k,
                         \widetilde{\mathcal M}_g\widehat\psi_k\rangle
       \notag\\
 &\le O(\log(k+1))+
       \sum_jw_j\log\langle\mathrm{post}_k,
                              \widetilde A_{j,g}\mathrm{post}_k\rangle.
 \label{comparator:sharp-jensen}
\end{align}
The conditioning cost is only $-\log q_k=O(\log(k+1))$.
For a leaf expectation in \eqref{comparator:sharp-jensen}, use
\eqref{comparator:pinned-metric} and scalar Hölder with powers three
on its three commuting signed label exponentials.  Each corresponding
post marginal is bounded by $z^{-k}\rho_Q^{\otimes k}$, by
\eqref{comparator:post-marginal}.  Apply
\eqref{comparator:signed-label-moments} at $3a$.  The result is
\[
 \log\langle\mathrm{post}_k,
                 \widetilde A_{j,g}\mathrm{post}_k\rangle
 \le ka g_{j,g}^{\mathrm{post}}+
                Cka^2\mathcal B+k|\log z|+O(\log(k+1)).
\]
Together with the pin norm this proves
\eqref{comparator:sharp-upper}.

\subsection{The inverse metric on copies near the ground state}

We next prove the lower comparison.  The gap first bounds the vector's
mass outside a subspace with few excited physical copies.  On that
subspace we will bound the compression of each inverse metric; this
will give an operator lower bound that can pass through the mean tree.
Let
$Q_i=\id-|\widetilde\Omega\rangle\langle\widetilde\Omega|$ on physical
copy $i$.  On the symmetric space with the two fixed auxiliary labels,
let $\Pi$ project onto $\sum_i Q_i\le\tau k$.  The full physical gap gives
\begin{equation}
 \overline{\widetilde H}-\widetilde E_0\id
       \ge\frac{\widetilde g}{k}\sum_iQ_i,
 \qquad
 \langle v,\Pi v\rangle\ge1-E_{\mathrm{def}}/\tau.
 \label{comparator:defect-mass}
\end{equation}
The projection is compatible with symmetry and the auxiliary sectors
because the defect count acts on physical copies only and is invariant
under simultaneous copy permutations.

Fix a unit vector $u$ in the range of $\Pi$.  For a subset
$B\subseteq\{1,\ldots,k\}$, let $u_B$ be its component with precisely
the physical copies in $B$ excited.  Then
\[
 u=\sum_{|B|\le\tau k}u_B,\qquad
 \sum_B\|u_B\|^2=1.
\]
In estimates for normalized components, omit the terms with $u_B=0$.
For $r=|B|$, $u_B$ is invariant under simultaneous permutations separately
on the $k'=k-r$ good copies and the $r$ bad copies.  It factors as
$\widetilde\Omega^{\otimes k'}$ on the good physical factors, tensored
with a vector on the remaining factors.  Each $u_B$ still has the two
whole-copy auxiliary labels $\lambda_k$.

We first compare metrics on the full copy space, before evaluating on
these components.  Extend a band metric there by its disjoint $W$ formula.
For any region $Q$, its whole label and its two subgroup labels commute.
On a compatible triple,
\begin{equation}
 F_Q^{\mathrm{good}}+F_Q^{\mathrm{bad}}
 \le F_Q^{\mathrm{whole}}
 \le F_Q^{\mathrm{good}}+F_Q^{\mathrm{bad}}+\log\binom{k}{r}.
 \label{comparator:restriction-dimensions}
\end{equation}
The first inequality is inclusion of a subgroup irreducible; the second
holds because its coset translates span the whole irreducible.  These are
the subgroup dimension inequalities from Lemma~\ref{lem:schur}.
For the three disjoint regions $P,Y,F$, all labels needed in this
comparison commute.  With
$G=F_P+F_F-F_Y$, the logarithmic approximation of the $W$ metrics yields
the operator inequality
\begin{equation}
 A_{j,g}^{-1}\le
 \poly(k)\binom{k}{r}^{Ca}
             \exp\{-aG^{\mathrm{good}}-aG^{\mathrm{bad}}\}.
 \label{comparator:whole-inverse}
\end{equation}
This step does not commute an individual physical-position projector
through $A_{j,g}$.  Such a commutation would generally be false.

On the bad group, simultaneous symmetry implies
$G^{\mathrm{bad}}\ge0$ on the measured support.  It may be discarded in
the expectation of \eqref{comparator:whole-inverse}; the good factors
preserve that support because they act on disjoint copies.  For each
auxiliary system, \eqref{comparator:restriction-dimensions} and its whole
label give, on $u_B$,
\begin{equation}
 F_C^{\mathrm{good}},\ F_R^{\mathrm{good}}
       \ge kS'-r\log d_*-\log\binom{k}{r}-o(k).
 \label{comparator:good-auxiliary}
\end{equation}
Indeed a bad-group Specht dimension is at most $d_*^r$, and
\eqref{comparator:high-label} controls the whole label.

The good copies therefore retain both large auxiliary labels.  To bound
the inverse metric on each component, we now compare the combined
physical--auxiliary labels with their separate labels.  The physical
factors will supply exact iid moments, and the loss from combining
labels will have exponential moments bounded by polynomials in $k$.
Suppress the good-group superscript for the next calculation and write
$Q=XU_{j,g}$, $V=V_{j,g}$, $Y=Y_{j,g}$.  Define the merge deficits
\[
 D_C=F_Q+F_C-F_{QC},\qquad
 D_R=F_V+F_R-F_{VR},\qquad
 G_{\mathrm{phys}}=F_Q+F_V-F_Y.
\]
All these operators, as well as $F_C,F_R$, commute.  The deficits are
nonnegative by the tensor-product dimension inequalities, and
\begin{equation}
 G^{\mathrm{good}}
       =G_{\mathrm{phys}}+F_C+F_R-D_C-D_R.
 \label{comparator:merge-decomposition}
\end{equation}
The physical good copies are symmetric, so they lie in the spectral
projection
$\mathsf J=\mathbf1_{[0,\infty)}(G_{\mathrm{phys}})$.
The projection $\mathsf J$ commutes with both merge deficits: physical
central labels commute with the larger-region simultaneous permutations
which define the merged labels.  Consequently it continues to enforce
$G_{\mathrm{phys}}\ge0$ while the other commuting factors are measured.
We use this compatible-physical-label projection, not an assertion that
the full physical symmetric subspace is preserved by a merge.

We need dimension-independent exponential rates for the deficits.
In a normalized component $u_B$, the marginal on $Q,C$ is the tensor
product of the iid physical marginal on $Q$ and a permutation-invariant
auxiliary marginal on $C$.  Conditional on their labels $\gamma,\beta$,
the two Specht factors are independent and maximally mixed.
If $\nu$ is their combined label, its conditional probability is
\[
       m_{\gamma,\beta}^{\nu}
             \frac{d_\nu}{d_\gamma d_\beta}.
\]
Here the multiplicity $m_{\gamma,\beta}^{\nu}$ is polynomial in $k$ at
fixed dimensions.  To see this, embed one copy of
$[\gamma]\otimes[\beta]$ into $(QC)^{\otimes k'}$ by choosing vectors
in the separate multiplicity spaces.  Its multiplicity of $[\nu]$
cannot exceed the dimension of the combined multiplicity space, which
is polynomial by Schur--Weyl decomposition.  Thus, for $0\le b\le1$,
\begin{equation}
 \E_{u_B/\|u_B\|}e^{bD_C}\le\poly(k),
 \qquad
 \E_{u_B/\|u_B\|}e^{bD_R}\le\poly(k).
 \label{comparator:merge-moments}
\end{equation}
For example the conditional sum is
$\sum_\nu m_{\gamma,\beta}^{\nu}
(d_\nu/(d_\gamma d_\beta))^{1-b}\le\sum_\nu m_{\gamma,\beta}^{\nu}$.
No iid hypothesis on the auxiliary marginal is used.  Its permutation
invariance follows by tracing the other factors from the good-group
symmetric component.

For the rough estimate, use $\mathsf J$ to discard
$e^{-aG_{\mathrm{phys}}}$, and use Cauchy--Schwarz on
$e^{aD_C}e^{aD_R}$ with \eqref{comparator:merge-moments}.
For the sharp estimate, use scalar Hölder with powers five on these two
factors and the three factors
$e^{-aF_Q},e^{-aF_V},e^{aF_Y}$.  Their physical moments are evaluated in
the exact iid state $\widetilde\Omega^{\otimes k'}$; apply
\eqref{comparator:signed-label-moments}.  Choosing $a_0$ and $c$ small
ensures that all fivefold moments and the negative-sign doubling in that
estimate are allowed.  Combining with
\eqref{comparator:good-auxiliary} gives
\begin{align}
 \frac{\langle u_B,A_{j,g}^{-1}u_B\rangle}{\|u_B\|^2}
 \le\poly(k)\exp\bigg\{
  -2kaS'-(k-r)a g_{j,g}^{\mathrm{pre}}
  +2ar\log d_*+Ca\log\binom{k}{r}
  +Cka^2\mathcal B+o(k)\bigg\}.
 \label{comparator:component-inverse}
\end{align}
For the rough estimate omit both the $g^{\mathrm{pre}}$ and
$a^2\mathcal B$ terms.  The entropy
$g_{j,g}^{\mathrm{pre}}$ is nonnegative because physical purity gives
$S(Y)=S(QV)\le S(Q)+S(V)$.

There are at most $(k+1)e^{kh(\tau)}$ subsets with $r\le\tau k$, and
$\log\binom{k}{r}\le kh(\tau)$.  Cauchy--Schwarz for
$A_{j,g}^{-1/2}\sum_Bu_B$, followed by
\eqref{comparator:component-inverse}, proves
\begin{equation}
 \Pi A_{j,g}^{-1}\Pi
       \le\exp\{-kaL_{j,g}+o(k)\}\Pi,
 \quad
 L_{j,g}=2S'+(1-\tau)g_{j,g}^{\mathrm{pre}}
          -C\left(\tau\log d_*+\frac{h(\tau)}a+a\mathcal B\right).
 \label{comparator:inverse-compression}
\end{equation}
In the rough version set $g^{\mathrm{pre}}=0$ and omit $a\mathcal B$.
The factor counting subsets is precisely the source of $h(\tau)/a$.
Taking the operator norm of $A_{j,g}^{-1/2}\Pi$ shows that
\eqref{comparator:inverse-compression} is equivalent to the useful lower
pin
\begin{equation}
 A_{j,g}\ge\exp\{kaL_{j,g}-o(k)\}\Pi.
 \label{comparator:lower-pin}
\end{equation}
For instance, apply the norm bound for $\Pi A_{j,g}^{-1/2}$ to
$A_{j,g}^{1/2}x$ to get the quadratic-form inequality for every $x$.
No commutation of $A_{j,g}$ with $\Pi$ is required.

\subsection{Passing the lower pin through the tree}

The individual lower pins must be combined with the floor on the whole
symmetric space before logarithms are taken.  For a common polynomial
$p(k)$, all the finitely many leaves satisfy
$A_{j,g}\ge p(k)^{-1}\id$ there.  This follows from
$F_Y\le F_P+F_F$ on simultaneous symmetry and the logarithmic metric
approximation.  Restrict throughout to the common auxiliary sector and
put $b_0=(2p(k))^{-1}$.  Averaging the floor and
\eqref{comparator:lower-pin} yields
\[
 A_{j,g}\ge b_0(\id-\Pi)+b_{j,g}\Pi,
 \qquad b_{j,g}=b_0+\tfrac12\exp\{kaL_{j,g}-o(k)\}.
\]
The right sides all commute with each other.  Mean monotonicity and
scalar homogeneity therefore give
\[
 \mathcal M_g\ge b_0(\id-\Pi)+e^{\ell_g}\Pi,
 \qquad \ell_g=\sum_jw_j\log b_{j,g}.
\]
Uniformly in the probabilities,
\[
 \ell_g\ge ka\sum_jw_jL_{j,g}-o(k),\qquad
 0\le\ell_g-\log b_0
       \le ka\max_j(L_{j,g})_++o(k).
\]
Log monotonicity and \eqref{comparator:defect-mass} now imply
\begin{equation}
 \langle v,\log\mathcal M_g\,v\rangle
 \ge ka\left[\sum_jw_jL_{j,g}
       -\frac{E_{\mathrm{def}}}{\tau}\max_j(L_{j,g})_+\right]-o(k).
 \label{comparator:tree-log-lower}
\end{equation}
One may first use $\min(1,E_{\mathrm{def}}/\tau)$ in the error, which
also makes its uniformity in the state immediate, and then weaken to the
displayed expression.  In the rough case
$\max_j(L_{j,g})_+\le2S'$; in the sharp case it is at most
$2S'+g_{\max}$.

Finally,
\[
 N^2\langle v,\mathcal M^{2s}v\rangle=\|\mathrm{pre}_k\|^2\le1,
 \qquad
 -\log N^2\ge\log\langle v,\mathcal M^{2s}v\rangle
       \ge2s\langle v,\log\mathcal M\,v\rangle.
\]
Since $\log\mathcal M=\sum_g\log\mathcal M_g$, summing
\eqref{comparator:tree-log-lower} proves
\eqref{comparator:rough-lower} and \eqref{comparator:sharp-lower}.

To finish, physical purity gives
$S(XV)=S(UY)$ and $S(XYV)=S(U)$.  Expanding the two mutual informations
on the right of \eqref{comparator:entropy-identity} proves that identity.
The comparison is therefore complete: its rough form discounts the target
entropy, while its sharp form bounds the target's mutual information with
the physical far systems.

All estimates above use maxima over finitely many statuses, fixed-dimension
Schur polynomial bounds, and the single sequence of labels
\eqref{comparator:high-label}.  The norm estimates for the defect
components are uniform in their auxiliary states.  Jensen and mean
monotonicity introduce no inverse terminal probability.  Consequently the
replica errors are uniform in all terminal probabilities, including at
zero weights by omitting those branches.  This also permits a later choice
of a different tree interpolation point for each $k$, followed by the
limit $k\to\infty$ at the fixed finite physical system.
